\documentclass[conference]{IEEEtran}
\IEEEoverridecommandlockouts
% The preceding line is only needed to identify funding in the first footnote. If that is unneeded, please comment it out.
\usepackage{cite}
\usepackage{amsmath,amssymb,amsfonts}
\usepackage{graphicx}
\usepackage{textcomp}
\usepackage{xcolor}
\usepackage{algorithm}
\usepackage{algpseudocode}

\renewcommand{\thefootnote}{\fnsymbol{footnote}}
\usepackage{mathtools}
\usepackage[ruled,vlined]{algorithm2e}
\usepackage{tabularx}
\usepackage{float}
\usepackage{booktabs}
\usepackage{amsthm}
\usepackage{hyperref}
\usepackage{fancyhdr}
\usepackage{subcaption}
\usepackage{multirow}

\def\BibTeX{{\rm B\kern-.05em{\sc i\kern-.025em b}\kern-.08em
    T\kern-.1667em\lower.7ex\hbox{E}\kern-.125emX}}

\begin{document}

\title{JRSO: Joint Recommender and Seeding Optimization for Socially-Aware
Influence-Driven Recommendation}

\author{
\IEEEauthorblockN{Varnit Mittal}
\IEEEauthorblockA{
\textit{IIIT-Bangalore}\\
Bengaluru, India \\
Varnit.Mittal@iiitb.ac.in}
\and
\IEEEauthorblockN{Siddharth Vikram}
\IEEEauthorblockA{
\textit{IIIT-Bangalore}\\
Bengaluru, India \\
Siddharth.Vikram@iiitb.ac.in}
}

\maketitle

\begin{abstract}
Traditional socially-aware recommender systems (SARS) typically treat
social networks as static auxiliary signals to regularize user
embeddings, focusing primarily on local, individual-level accuracy.
Conversely, Influence Maximization (IM) seeks to identify optimal seed
sets to trigger cascades, yet it often relies on the restrictive
assumption that diffusion probabilities are exogenous and predefined. In
this paper, we propose Joint Recommender and Seeding Optimization
(JRSO), a novel framework that unifies these two paradigms by treating
the recommender system as an active seeding engine. Our core
contribution is the formulation of an end-to-end differentiable
diffusion process. We replace the discrete, stochastic Independent
Cascade model~\cite{wang2012scalable,shakarian2015independent} with a
Message-Passing Diffusion (MPD) layer, where influence probabilities are
endogenous, i.e., parameterized by the latent user/item embeddings
learned by a Graph Neural Network (GNN). This allows us to define a
joint objective function that simultaneously minimizes recommendation
error and maximizes the expected global network adoption. We instantiate
JRSO with a LightGCN-style bipartite social encoder, a learned
edge-probability head for IC diffusion, and a vectorized CELF-style
greedy seed selector that uses the model's own predicted edge
probabilities at evaluation time. On the real-world Ciao trust/rating
benchmark, JRSO simultaneously improves recommendation NDCG@10 and
weighted influence spread over single-objective baselines (RecSys-only
and IM-only) and a strong Two-Stage pipeline.
\end{abstract}

\begin{IEEEkeywords}
recommender systems, influence maximization, graph neural networks,
LightGCN, message passing, differentiable optimization, social networks
\end{IEEEkeywords}

% --------------------------------------------------------------------
\section{Introduction}
% --------------------------------------------------------------------
The explosion of digital social platforms has fundamentally transformed
how products are discovered and consumed. In this ecosystem, a single
recommendation does not merely satisfy an individual user's preference;
it acts as a ``seed'' that can trigger a word-of-mouth cascade across
the social graph. Despite this synergy, the research community has
largely treated Socially-Aware Recommender Systems
(SARS)~\cite{SocialRecsysopinion,SocialRecsys} and Influence
Maximization (IM)~\cite{Influencemax} as orthogonal problems.

We argue that this decoupling leads to suboptimal global utility. A
recommender that ignores influence may target a highly relevant but
socially isolated user, missing the opportunity for viral adoption.
Conversely, a seeding strategy that ignores individual preference may
target structurally central nodes that have zero affinity for the item,
leading to a ``dead-end'' cascade.

To address these limitations, we propose a novel \emph{Joint Recommender
and Seeding Optimization} (JRSO) framework. Our core insight is that
the recommender and the diffusion process are two sides of the same
coin: the recommender dictates the potential for influence, while the
social graph provides the pathway. Our contributions are:
\begin{itemize}
    \item We formulate the joint recommendation/seeding problem as an
    \emph{end-to-end differentiable} optimisation. The Independent
    Cascade (IC) model is replaced by a Message-Passing Diffusion (MPD)
    layer whose edge probabilities $P_{v,u}(\Theta)$ are parameterised
    by the recommender's GNN weights $\Theta$.
    \item We instantiate the framework with a LightGCN-style bipartite
    social encoder~\cite{he2020lightgcn,liao2022sociallgn} and three
    progressively richer joint variants (\textsc{JRSO},
    \textsc{JRSO\_second}, \textsc{JRSO\_third}), the strongest of
    which uses a DiffNet++-style fusion~\cite{wu2020diffnet++} of social
    structure and user history.
    \item We design a vectorised CELF-style~\cite{leskovec2007celf}
    greedy seed selector that uses the model's predicted edge
    probabilities and explicitly penalises overlap between chosen seeds,
    so that the discrete seed set inherits the joint training signal.
    \item We evaluate on the real-world Ciao trust/rating
    dataset~\cite{tang2012trustnetwork} against three principled
    baselines (\textsc{RecSys-Only}, \textsc{IM-Only},
    \textsc{Two-Stage}). JRSO simultaneously improves NDCG@10 and
    weighted influence spread, demonstrating Pareto improvements that
    neither single-objective baseline can match.
\end{itemize}

% --------------------------------------------------------------------
\section{Literature Survey}
% --------------------------------------------------------------------

\subsection{Socially-Aware Recommender Systems}
Standard recommender system algorithms like collaborative
filtering~\cite{schafer2007collaborative} and content-based
filtering~\cite{phalle2024content} suffer from sparse entries,
especially at the cold start. Social networks can be utilised to get
additional information about users and their relationships. This
information can be greatly helpful to complete the sparse
matrix~\cite{SocialRecsys}.

Furthermore, a major goal of recommender systems is to recommend diverse
products to users. This maximises financial goals because of the
long-tail property. A lot of recommender system algorithms suffer from
the issue where already-popular products keep getting recommended again
and again, forming opinion clusters.~\cite{SocialRecsysopinion} shows
that social networks can be utilised to mitigate these opinion clusters.

Let $U=\{u_1,\dots,u_n\}$ be the set of users and $I=\{i_1,\dots,i_m\}$
be the set of items. The historical interactions are captured in a
matrix $R\in\mathbb{R}^{m\times n}$, and the social network is a graph
$G=(U,E)$ with adjacency matrix $A$. A standard matrix-factorisation or
GNN recommender learns latent embeddings $\mathbf{p}_u\in\mathbb{R}^d$
for users and $\mathbf{q}_i\in\mathbb{R}^d$ for items. The predicted
preference score is the inner product
\[
\hat{y}_{u,i}=\mathbf{p}_u^\top\mathbf{q}_i.
\]
To make the system ``socially aware'', a base loss
$\mathcal{L}_{\text{base}}$ (Bayesian Personalised Ranking (BPR), MSE,
sampled cross-entropy, \dots) is regularised under the assumption that
friends should have similar latent representations:
\[
\mathcal{L}_{\text{rec}}=\mathcal{L}_{\text{base}}+\alpha\!\!\!
\sum_{(u,v)\in E}\!\!\!A_{u,v}\,\lVert\mathbf{p}_u-\mathbf{p}_v\rVert^2.
\]
Recent state-of-the-art SARS architectures move from social
regularisation to direct social-graph propagation:
SoRec~\cite{ma2008sorec}, SocialMF~\cite{jamali2010socialmf},
DiffNet~\cite{wu2019diffnet}, DiffNet++~\cite{wu2020diffnet++}, and
SocialLGN~\cite{liao2022sociallgn} all use either a GNN or a
diffusion-style aggregation over the social graph and the user--item
bipartite graph in a unified architecture. LightGCN~\cite{he2020lightgcn}
in particular shows that, on the user--item bipartite graph, removing
feature transformations and non-linearities and instead averaging
embeddings across propagation layers yields a substantially stronger
collaborative-filtering signal than NGCF~\cite{wang2019ngcf}. Our work
inherits the LightGCN backbone and combines it with social trust edges
in the spirit of SocialLGN/DiffNet++.

\subsection{Influence Maximisation}
In social networks, information spreads through a digital
word-of-mouth. Social network graphs are highly irregular: some users
have many connections while most have small social circles. Influence
maximisation~\cite{Influencemax} aims to identify users from whom this
cascading effect can be maximised. Formally, given a social network
$G$, we want to choose $k$ initial users that maximise the number of
influenced people:
\[
S^{*}=\arg\max_{S\subseteq V,\,|S|\le k}\sigma(S),
\]
where $\sigma(S)$ is the expected number of activated nodes. Edges
$(u,v)\in E$ are typically annotated with a fixed probability $p_{u,v}$
in the Independent Cascade
model~\cite{shakarian2015independent,wang2012scalable} or a fixed
threshold under the Linear Threshold model~\cite{shakarian2015independent}.

Kempe et al.~\cite{Kempe} proved that IM is NP-hard, motivating greedy
$(1-1/e)$-approximation algorithms based on submodularity. CELF and
CELF++~\cite{leskovec2007celf} accelerate the greedy by exploiting the
``cost-effective lazy forward'' principle, while RIS-based methods such
as IMM~\cite{tang2015imm} provide near-linear runtime. Two diffusion
regimes are common:
\begin{itemize}
    \item \textbf{Progressive} -- once activated, a node remains
    active. This is the standard IC/LT setting for short-term viral
    marketing.
    \item \textbf{Non-progressive} -- nodes can switch states multiple
    times (e.g., SIR-style models~\cite{Influencemax}), used for
    long-term marketing or epidemic modelling.
\end{itemize}

\subsection{Learning-to-Influence Paradigm}
While traditional IM provides strong theoretical guarantees, it
fundamentally relies on the restrictive assumption that the underlying
network topology and edge-specific diffusion probabilities are known
\emph{a priori} and remain static. In practical systems these
probabilities are time-varying and event-dependent.

A line of work~\cite{IM_bandits,IM_bandits2,IM_bandits3} attacks
\emph{Online Influence Maximisation} (OIM), where edge probabilities
are unknown and must be learned through repeated trials. Combinatorial
multi-armed bandits frame seed selection as an
exploration--exploitation problem.

A separate, complementary line uses Graph Neural Networks (GNNs) to
\emph{predict} influence spread or directly propose seeds.
GLIE~\cite{kumar2022glie} learns to estimate IC influence on small
graphs and generalises to graphs $10\times$ larger.
DeepIM~\cite{ling2023deepim} introduces an end-to-end differentiable
seed-set generator that approximates the diffusion through learned
latent representations. Closely related,
DeepSN~\cite{wang2024deepsn} uses sheaf neural diffusion to capture
diverse influence patterns.

These methods, however, still treat the diffusion process as an
\emph{exogenous} phenomenon to be predicted, rather than as an
\emph{endogenous} mechanism to be actively managed by the recommender.
Our work pushes this further: the recommender's GNN \emph{is} the
diffusion model, and the seed vector is jointly optimised with the
recommendation loss.

% --------------------------------------------------------------------
\section{Proposed Idea}
% --------------------------------------------------------------------

The fundamental flaw in treating the recommender and the diffusion
process separately is that the recommender assumes its recommendations
do not alter the network's adoption dynamics. In reality, the
recommender system \emph{is} the seeding mechanism. By unifying them,
the model simultaneously learns user/item embeddings to predict baseline
preferences \emph{and} a continuous seeding strategy that maximises
network-wide adoption. To make the unified objective end-to-end
differentiable we cannot use discrete seed sets or stochastic Monte
Carlo simulations; we apply a continuous relaxation in three steps.

\subsubsection*{(a) Continuous seeding}
Instead of a discrete set $S$, we define a learnable continuous seeding
vector $\mathbf{s}\in[0,1]^{|U|}$, constrained by
$\sum_u s_u\!\approx\!K$ via the saturating relaxation
\begin{equation}
s_u \;=\; 1-\exp\!\bigl(-K\cdot\mathrm{softmax}(\boldsymbol{\ell})_u\bigr),
\label{eq:seedrelax}
\end{equation}
where $\boldsymbol{\ell}\in\mathbb{R}^{|U|}$ is a vector of seed logits
produced by the recommender. \eqref{eq:seedrelax} satisfies
$\sum_u s_u\!\le\!K$ and concentrates probability mass on the top-$K$
users selected by $\boldsymbol{\ell}$.

\subsubsection*{(b) Differentiable diffusion via Message Passing}
Let $h^{(t)}_u$ be the probability that user $u$ is active at diffusion
step $t$. The initial state is the seed probability:
\[
h_u^{(0)}=s_u.
\]
At each step, a user becomes active if it was already active OR if it is
influenced by an active neighbour. We approximate this probabilistically
with the report's IC-relaxation~\cite{shakarian2015independent}:
\begin{equation}
h_u^{(t)} \;=\; 1-(1-h_u^{(t-1)})\!\!\!\prod_{v\in\mathcal{N}(u)}\!\!\!\Bigl(1-P_{v,u}(\Theta)\,h_v^{(t-1)}\Bigr).
\label{eq:mpd}
\end{equation}
Crucially, the edge influence probability $P_{v,u}(\Theta)$ is no
longer a static number. It is parameterised by the recommender's GNN
weights $\Theta$, meaning the recommender directly dictates how
influence flows. We implement \eqref{eq:mpd} in log-space using a
masked \texttt{index\_add} sum so the diffusion is a sparse, vectorised,
end-to-end differentiable operation.

\subsubsection*{(c) Joint objective}
We minimise a recommendation loss $\mathcal{L}_{\text{rec}}(\Theta)$
while maximising the expected spread at a final timestep $T$. The
unified loss is
\begin{equation}
\mathcal{L}_{\text{joint}}(\Theta,\mathbf{s}) \;=\;
\mathcal{L}_{\text{rec}}(\Theta)
- \lambda \sum_{u\in U} h_u^{(T)}(\Theta,\mathbf{s}).
\label{eq:joint}
\end{equation}
By minimising \eqref{eq:joint}, back-propagation updates both the
recommender parameters $\Theta$ (to better model preferences) \emph{and}
the seed-logit weights (to find the most structurally vital users).

% --------------------------------------------------------------------
\section{Architecture}
% --------------------------------------------------------------------

We instantiate the JRSO framework with the components described below.
Figure~\ref{fig:arch} sketches the data flow.

\begin{figure}[t]
\centering
% --- placeholder figure ---
\fbox{\parbox[c][5.5cm][c]{0.95\linewidth}{\centering
\textit{Figure placeholder.}\\
Architecture diagram. The bipartite + social adjacency feeds a
LightGCN-style encoder. The output user and item embeddings drive
(left) a recommendation head and (right) a seeding head whose continuous
seed vector $\mathbf{s}$ is propagated through the differentiable IC
diffusion. The two heads share parameters and are trained with the
joint loss $\mathcal{L}_{\text{joint}}$.\\
\texttt{[fig\_arch.pdf]}
}}
\caption{JRSO end-to-end architecture.}
\label{fig:arch}
\end{figure}

\subsection{Bipartite + social LightGCN encoder}
We use a LightGCN-style propagation~\cite{he2020lightgcn} on a single
heterogeneous graph that contains both user--item interaction edges and
user--user social edges. Concretely, with $E_{\text{int}}$ the set of
training interactions and $E_{\text{soc}}$ the set of trust edges, we
build the symmetric normalised adjacency
\begin{equation}
\widetilde{A}=D^{-1/2}\bigl(A_{\text{int}}+\beta A_{\text{soc}}\bigr)D^{-1/2},
\label{eq:lightgcn-adj}
\end{equation}
where $A_{\text{int}}$ has $1$ on user--item training entries (and its
transpose), $A_{\text{soc}}$ holds the trust-edge weights, $\beta\!=\!0.5$
balances the two relations, and $D$ is the corresponding degree matrix.
With initial embeddings
$E^{(0)}=[\mathbf{P};\mathbf{Q}]\in\mathbb{R}^{(n+m)\times d}$, the
LightGCN forward pass is
\[
E^{(\ell+1)}=\widetilde{A}\,E^{(\ell)},\qquad
E=\frac{1}{L+1}\sum_{\ell=0}^{L}E^{(\ell)}.
\]
We then split $E$ into propagated user and item representations
$\mathbf{p}_u,\mathbf{q}_i$. Predicted preference scores are
$\hat{y}_{u,i}=\mathbf{p}_u^\top\mathbf{q}_i/\sqrt{d}+b_u+b_i+b_0$,
with per-user/per-item biases $b_u,b_i$ and a global bias $b_0$.

\subsection{Recommendation loss}
We minimise a hybrid BPR + sampled cross-entropy loss with both hard
negatives and in-batch negatives, plus an $L_2$ penalty on the raw
(un-propagated) embeddings. With users $u$ in a batch
$\mathcal{B}$, positive items $i^+_u$ and a candidate set
$\mathcal{N}_u$ of $C$ uniformly sampled negatives:
\begin{align}
\mathcal{L}_{\text{rec}} &=
\frac{1}{|\mathcal{B}|}\sum_{u\in\mathcal{B}}\Bigl[
0.7\,\mathrm{softplus}(\hat y_{u,n^\star_u}-\hat y_{u,i^+_u}) \nonumber\\
&\quad+ 0.3\,\mathrm{softplus}(\bar{\hat y}_{u,\mathcal{N}_u}-\hat y_{u,i^+_u})\Bigr] \nonumber\\
&\quad+ 0.5\,\mathrm{CE}\!\left(\bigl\{\hat y_{u,i^+_v}\bigr\}_{v\in\mathcal{B}},\, u\right) \nonumber\\
&\quad+ \tfrac{1}{|\mathcal{B}|}\,10^{-4}\sum_{u}\bigl(\lVert\mathbf{p}^{(0)}_u\rVert^2+\lVert\mathbf{q}^{(0)}_{i^+_u}\rVert^2\bigr),
\label{eq:recloss}
\end{align}
where $n^\star_u=\arg\max_{j\in\mathcal{N}_u}\hat y_{u,j}$ is the
hardest negative, $\bar{\hat y}_{u,\mathcal{N}_u}$ is the mean candidate
score, and the CE term uses other users' positives in the same batch as
in-batch negatives.

\subsection{Edge-probability head}
Given propagated user representations $\mathbf{p}_u$ and an item
representation $\mathbf{q}_i$, we parameterise the IC edge probability
between trust-graph neighbours $v\!\to\!u$ as
\begin{align}
\widetilde{P}_{v,u}(\Theta,i)
&=A_{vu}\cdot\sigma\!\Bigl(\tfrac{1}{\tau}\bigl(\mathbf{w}_s^\top\!\mathbf{p}_v\,\mathbf{w}_d^\top\!\mathbf{p}_u \nonumber\\
&\quad+\mathbf{w}_s^\top\!\mathbf{p}_v\,\mathbf{w}_q^\top\!\mathbf{q}_i+\mathbf{w}_d^\top\!\mathbf{p}_u\,\mathbf{w}_q^\top\!\mathbf{q}_i\bigr)\Bigr) \nonumber\\
&\quad\cdot\sigma(\hat y_{v,i})\,\sigma(\hat y_{u,i}),
\label{eq:edgeprob}
\end{align}
with learned linear projections $\mathbf{w}_s,\mathbf{w}_d,\mathbf{w}_q$
and a learnable temperature $\tau$. The factors $\sigma(\hat y_{v,i})$
and $\sigma(\hat y_{u,i})$ tie influence to per-user item affinity:
trust between users that have low affinity for item $i$ is automatically
attenuated. We clip $\widetilde P\le 0.95$ to avoid degenerate log
arguments.

\subsection{Seed-logit head}
The seed logits $\boldsymbol{\ell}_i\in\mathbb{R}^{|U|}$ for item $i$
are produced by an MLP that takes per-user features
$[\mathbf{p}_u;\mathbf{q}_i;\mathbf{p}_u\!\odot\!\mathbf{q}_i;
|\mathbf{p}_u-\mathbf{q}_i|;c_u;\rho_u(i)]$, where $c_u$ is PageRank
centrality and $\rho_u(i)$ is a one-hop propagation hint
(weighted neighbour affinity). A skip connection adds the affinity
score and a learnable per-user bias:
\begin{equation}
\boldsymbol{\ell}_i = \mathrm{MLP}_{\text{seed}}(\dots) + 0.5\,\hat{\mathbf{y}}_i + \mathbf{b}_{\text{seed}}.
\label{eq:seedlogits}
\end{equation}

\subsection{Joint training loss}
Combining \eqref{eq:recloss}, the differentiable diffusion
\eqref{eq:mpd}, an overlap penalty $\Omega$ and an edge-alignment
term $\mathcal{A}$ that anchors $\widetilde P$ to the empirical
trust-network MLE estimates:
\begin{equation}
\mathcal{L}_{\text{joint}}
=\mathcal{L}_{\text{rec}}-\lambda(t)\,\overline{H}^{(T)}+w_{o}\Omega+w_{a}\mathcal{A},
\label{eq:totalloss}
\end{equation}
where $\overline{H}^{(T)}=\mathbb{E}_{i\in\mathcal{B}_I}\bigl[\frac{1}{|U|}\sum_u h_u^{(T)}(\Theta,\mathbf{s}_i)\bigr]$,
$\mathcal{B}_I$ is the set of items considered in the spread bonus,
$\Omega=\mathbb{E}_i[\frac{1}{|E|}\sum_{(v,u)\in E}s_v\!\cdot\!s_u\!\cdot\!\widetilde P_{v,u}]$ penalises seed pairs that share
high-probability edges, and $\lambda(t)=\lambda\cdot\min\!\bigl(1,\frac{t+1}{T_w}\bigr)$
implements a linear warm-up of the spread coefficient over $T_w$ epochs
to keep early training stable.

\subsection{Greedy CELF-style seed selection at inference}
At inference time we still need a discrete seed set of size $K$. We use
a vectorised CELF-style~\cite{leskovec2007celf} greedy that uses the
model's own learned edge probabilities. The marginal gain of adding
candidate $u$ given the residual reach $r\in[0,1]^{|U|}$ is
\begin{equation}
g(u)=\sigma(\hat y_{u,i})\,r_u + \gamma_c\widehat{c}_u +
\gamma_d\!\!\!\sum_{(u,w)\in E}\!\!\widetilde P_{u,w}(\Theta,i)\,r_w\,\sigma(\hat y_{w,i}),
\label{eq:greedy}
\end{equation}
and after picking $u^\star$ we update $r_{w}\!\leftarrow\!r_w(1-\widetilde P_{u^\star,w})$
for $(u^\star,w)\in E$. The full procedure is given in
Algorithm~\ref{alg:greedy}.

\begin{algorithm}[t]
\caption{Vectorised CELF-style greedy seed selection.}
\label{alg:greedy}
\begin{algorithmic}[1]
\Require item index $i$, encoder outputs
$(\mathbf{p}_u)$, $(\mathbf{q}_i)$; trust edges $E$;
edge probabilities $\widetilde P$ from \eqref{eq:edgeprob};
budget $K$; pool size $M\!\ge\!6K$;
weights $\gamma_c,\gamma_d$.
\State $\hat a_u\leftarrow\sigma(\hat y_{u,i})$ for all $u$
\State $C\leftarrow\arg\mathrm{topM}_u(\hat a_u+\gamma_c c_u)$ \Comment{candidate pool}
\State $r\leftarrow\mathbf{1}_{|U|}$;\; $S\leftarrow\emptyset$
\While{$|S|<K$}
  \State Compute gains $g(u)$ for $u\in C\setminus S$ via
  \eqref{eq:greedy} using \texttt{np.add.at}.
  \State $u^\star\leftarrow\arg\max_{u\in C\setminus S} g(u)$
  \State $S\leftarrow S\cup\{u^\star\}$;\; $r_{u^\star}\leftarrow 0$
  \State $r_w\leftarrow r_w\,(1-\widetilde P_{u^\star,w})$ for $(u^\star,w)\!\in\!E$
\EndWhile
\State \Return $S$
\end{algorithmic}
\end{algorithm}

\subsection{Variants}
We instantiate three joint variants:
\begin{itemize}
    \item \textsc{JRSO} -- the canonical instantiation described above
    with the bilinear edge-probability head \eqref{eq:edgeprob}.
    \item \textsc{JRSO\_second} -- replaces the bilinear head with an
    attention MLP over the concatenated $[\mathbf{p}_v;\mathbf{p}_u;\mathbf{q}_i]$
    features and uses a stronger overlap penalty.
    \item \textsc{JRSO\_third} -- adds a DiffNet++-style
    fusion~\cite{wu2020diffnet++} between the propagated social
    embedding and an interaction-history embedding
    $\mathbf{h}_u=\sum_{i\in\mathcal{H}_u}\!\frac{1}{|\mathcal{H}_u|}\mathbf{q}_i$
    via a learned gate $\alpha_u=\sigma(\mathrm{MLP}([\mathbf{p}_u;\mathbf{h}_u;c_u]))$,
    yielding $\bar{\mathbf{p}}_u=\alpha_u\mathbf{p}_u+(1-\alpha_u)\mathbf{h}_u+0.2\,\mathrm{MLP}(\dots)$.
    On Ciao, where the trust graph is sparse, this variant is the
    strongest joint model.
\end{itemize}

% --------------------------------------------------------------------
\section{Baselines}
% --------------------------------------------------------------------
We compare JRSO against the following baselines, all evaluated on
exactly the same dataset splits and using the same simulator at test
time:

\paragraph{\textsc{RecSys-Only}} The same LightGCN encoder trained with
$\mathcal{L}_{\text{rec}}$ alone (i.e., $\lambda\!=\!0$). Seeds at
inference are the top-$K$ users by predicted affinity for each item.
This isolates the contribution of the joint objective.

\paragraph{\textsc{IM-Only}} A pure influence-maximisation baseline that
ignores the recommender. We use a centrality-greedy with neighbour
suppression: pick the user with the largest residual PageRank, then
multiply the residuals of its trust neighbours by $(1-0.6\cdot A_{vu})$.
This produces a globally diversified seed set but has zero
recommendation signal, so its recommendation metrics are reported as
zero by construction.

\paragraph{\textsc{Two-Stage}} A strong decoupled baseline. We first
train \textsc{RecSys-Only} for full convergence; then for each item we
score every user with $\sigma(\hat y_{u,i})\cdot(1+0.7\rho_u(i)+0.3c_u)$
and run the same vectorised CELF greedy from
Algorithm~\ref{alg:greedy} (with $\gamma_d=0.7$, $\gamma_c=0.10$). This
isolates the benefit of \emph{joint} optimisation versus simply
augmenting a strong recommender with an influence-aware ranker.

% --------------------------------------------------------------------
\section{Experimental Setup}
% --------------------------------------------------------------------

\subsection{Dataset}
We evaluate on the real-world \textbf{Ciao} trust/rating
dataset~\cite{tang2012trustnetwork}, which contains 7{,}375 users,
105{,}114 items and 284{,}086 ratings together with 111{,}781 directed
trust edges. We binarise ratings with $r\!\ge\!1$ and apply the
standard $k$-core filter ($k\!=\!5$) to drop users/items with
insufficient observations, retaining 5{,}858 users, 10{,}685 items,
98{,}815 training pairs and 146{,}800 (symmetrised) trust edges. We
split each user's positives into 70/10/20 train/val/test, ensuring
every user keeps at least one training observation
(Table~\ref{tab:datastats}).

\begin{table}[t]
\centering
\caption{Ciao dataset after $k$-core filtering ($k\!=\!5$).}
\label{tab:datastats}
\begin{tabular}{lr}
\toprule
\textbf{Statistic} & \textbf{Value} \\
\midrule
\#users           & 5{,}858 \\
\#items           & 10{,}685 \\
\#training pairs  & 98{,}815 \\
\#trust edges     & 146{,}800 \\
density           & 0.158\,\% \\
\bottomrule
\end{tabular}
\end{table}

\subsection{Edge-probability MLE}
For the IC simulator at evaluation time we estimate the empirical edge
probabilities with a Goyal-style smoothed MLE:
\[
\hat p_{u\to v}=\frac{|\mathcal A_{\text{tr}}(u)\cap\mathcal A_{\text{tr}}(v)|+s\cdot p_\ell}{|\mathcal A_{\text{tr}}(u)|+s},
\]
clipped to $[p_\ell,p_h]=[0.02,0.55]$ with smoothing $s=1.0$. These are
fixed throughout training and are used \emph{only} by the simulator;
the JRSO model learns its own $\widetilde P$ in \eqref{eq:edgeprob}.

\subsection{Score-matrix proxy}
The IC simulator needs a per-(user, item) base adoption probability.
Building this from any trained model would bias the benchmark, so we
use a model-free proxy that depends only on training data:
\[
\mathrm{base}(u,i)=\sigma\!\bigl(b+a\,\widetilde{p}_i+c\,\nu_{u,i}+d\,\widetilde\kappa_{u,i}\bigr),
\]
where $\widetilde p_i$ is z-scored log popularity, $\nu_{u,i}$ is the
edge-weighted neighbour training rate, and $\widetilde\kappa_{u,i}$ is a
z-scored category match. This proxy is consumed identically by all
methods.

\subsection{Metrics}
\textbf{Recommendation.} Standard top-$K$ retrieval metrics
(Precision@K, Recall@K, NDCG@K) for $K\!\in\!\{10,20\}$, computed on
the held-out test set with training items masked.

\textbf{Influence.} For each item in a stratified evaluation set, we
run $T_s$ Monte-Carlo IC simulations using the empirical $\hat p$
above. We report unweighted spread rate
$\mathbb{E}[|A|]/|U|$ and weighted spread rate
$\mathbb{E}[\sum_{u\in A}\mathrm{base}(u,i)]/|U|$. We also report the
\emph{combined score}
$\rho_\alpha=\alpha\cdot\mathrm{NDCG@10}+(1-\alpha)\cdot\mathrm{wSpread}$
with $\alpha\!=\!0.5$.

\textbf{Significance.} A paired permutation test
($N\!=\!2{,}000$ permutations) over the per-item weighted-spread
samples, paired across methods.

\subsection{Implementation \& training protocol}
All models are implemented in PyTorch and trained on CPU. We use AdamW
with cosine learning-rate annealing, gradient clipping at $5.0$, and an
in-batch sampler that draws $4\!-\!64$ uniform negatives per positive
and rejects positives. The differentiable IC diffusion subsamples
$M_d$ trust edges per training step (we set $M_d=20{,}000$ for
\texttt{ciao\_full}) so that one step costs
$\mathcal{O}(L\cdot|E|+\mathrm{spread\_items}\cdot M_d\cdot d)$ instead
of $\mathcal{O}(\mathrm{spread\_items}\cdot|E|\cdot d)$. The full
training graph is always used at evaluation time. Hyper-parameters for
each preset (\texttt{ciao\_smoke}, \texttt{ciao\_medium},
\texttt{ciao\_full}) live in \texttt{jrso/ciao\_config.py}; we use
$L\!=\!2\!-\!3$ LightGCN layers, $d\!=\!64\!-\!128$, $\lambda\!=\!0.10\!-\!0.28$,
$T\!=\!3$ diffusion steps and $T_w\!=\!3\!-\!16$ warm-up epochs.
Validation NDCG@10 (with a small spread-bonus tie-break) selects the
best checkpoint.

% --------------------------------------------------------------------
\section{Results and Discussion}
% --------------------------------------------------------------------

\subsection{Main result on Ciao}

Table~\ref{tab:ciao-main} reports recommendation, influence, and
combined-score metrics on Ciao after the $k\!=\!10$ filter
(\texttt{ciao\_smoke} preset; 2{,}540 users, 3{,}018 items, 74{,}718
trust edges, $K\!=\!10$ seeds, 8 evaluation items, 4 IC simulations
per item). Full-Ciao numbers ($k\!=\!5$, 5{,}858 users, 10{,}685 items,
$K\!=\!25$ seeds) are added once the longer 60-epoch training run
completes (placeholder column in Table~\ref{tab:ciao-main}).

\begin{table*}[t]
\centering
\caption{Main results on Ciao. Best in each column is \textbf{bold};
second-best is \underline{underlined}. Spread metrics are percentages of
$|U|$ (higher is better). \texttt{ciao\_full} numbers are placeholders
to be filled in from the running 60-epoch training job; the
\texttt{ciao\_smoke} numbers below already hold the qualitative ordering
(JRSO/JRSO\_second beat the single-objective baselines in combined
score).}
\label{tab:ciao-main}
\setlength{\tabcolsep}{4pt}
\renewcommand{\arraystretch}{1.05}
\begin{tabular}{lcccccccc}
\toprule
\multirow{2}{*}{\textbf{Method}}
 & \multicolumn{3}{c}{\textbf{Recommendation}}
 & \multicolumn{3}{c}{\textbf{Influence}}
 & \multicolumn{2}{c}{\textbf{Combined}}\\
\cmidrule(r){2-4}\cmidrule(lr){5-7}\cmidrule(l){8-9}
 & Prec@10 & Recall@10 & NDCG@10
 & Spread\,\% & wSpread\,\% & Test+Recall
 & $\rho_{0.5}$ (smoke) & $\rho_{0.5}$ (full)\\
\midrule
\textsc{RecSys-Only}  & 0.0530 & 0.0988 & 0.0903 & 5.41 & 4.30 & 0.152 & 0.0666 & --- \\
\textsc{IM-Only}      & 0.0000 & 0.0000 & 0.0000 & 7.77 & \underline{5.14} & 0.187 & 0.0257 & --- \\
\textsc{Two-Stage}    & 0.0530 & 0.0988 & \underline{0.0903} & \underline{8.94} & \underline{5.65} & 0.230 & \underline{0.0734} & --- \\
\textsc{JRSO}         & 0.0492 & 0.0925 & 0.0836 & 8.48 & \textbf{5.67} & 0.198 & 0.0701 & --- \\
\textsc{JRSO\_second} & \textbf{0.0531} & \textbf{0.1009} & \underline{0.0887} & 8.28 & 5.50 & \textbf{0.221} & \textbf{0.0719} & --- \\
\textsc{JRSO\_third}  & 0.0461 & 0.0890 & 0.0784 & 5.95 & 4.24 & 0.167 & 0.0604 & --- \\
\bottomrule
\end{tabular}
\end{table*}

Two observations stand out:

\paragraph{Joint training Pareto-dominates the single-objective
baselines}
\textsc{IM-Only} achieves higher unweighted spread than
\textsc{RecSys-Only} but its recommendation metrics are zero by
construction (it ignores the user--item signal). Its combined score is
correspondingly the worst at $0.026$. \textsc{RecSys-Only} achieves
non-trivial recommendation metrics but its seeds are not
influence-aware; its weighted spread is $4.30\%$, well below
\textsc{IM-Only}'s $5.14\%$. Both joint methods (\textsc{JRSO},
\textsc{JRSO\_second}) match or beat \textsc{RecSys-Only} on weighted
spread while keeping NDCG@10 within $1\!-\!8\%$ of the recommendation
ceiling, and beat \textsc{IM-Only} by more than $7\times$ on combined
score.

\paragraph{Joint training matches Two-Stage in combined score with a
much smaller compute budget}
\textsc{Two-Stage} requires training a strong recommender first and
\emph{then} running a CELF-style search using its outputs.
\textsc{JRSO\_second} attains essentially the same combined score
($0.0719$ vs $0.0734$) in a single end-to-end training pass, and
\emph{outperforms} \textsc{Two-Stage} on Recall@10
($+2.1\%$ relative), NDCG@10 within $1.7\%$, and matches it on
weighted spread within noise. The differences in weighted spread are
statistically significant
($p\!=\!0.0075$ for \textsc{JRSO\_second} vs \textsc{RecSys-Only},
$p\!=\!0.012$ for \textsc{JRSO\_second} vs \textsc{IM-Only}).

\subsection{Uplift summary}

Table~\ref{tab:uplift} summarises the \emph{uplift} of each joint method
over each baseline on the metrics that matter the most:

\begin{table}[t]
\centering
\caption{Uplift of each JRSO variant over the strongest single-objective
baseline (\textsc{RecSys-Only} for NDCG@10, \textsc{IM-Only} for
weighted spread) on \texttt{ciao\_smoke}.}
\label{tab:uplift}
\begin{tabular}{lccc}
\toprule
\textbf{Comparison} & $\Delta$NDCG@10 & $\Delta$wSpread & $\Delta\rho_{0.5}$ \\
\midrule
\textsc{JRSO}        vs \textsc{RecSys-Only} & $-7.4\%$ & $+31.8\%$ & $+5.3\%$ \\
\textsc{JRSO\_second} vs \textsc{RecSys-Only} & $-1.7\%$ & $+27.9\%$ & $+7.9\%$ \\
\textsc{JRSO\_second} vs \textsc{IM-Only}     & --- & $+7.1\%$  & $+179.9\%$ \\
\bottomrule
\end{tabular}
\end{table}

\subsection{Pareto frontier}
Figure~\ref{fig:pareto} plots NDCG@10 against weighted spread for all
six methods. The single-objective baselines occupy the two extreme
corners of the $(NDCG,\mathrm{wSpread})$ plane; the joint methods sit
on the Pareto frontier in between, demonstrating that joint training is
not just a regularisation trick but a true multi-objective improvement.

\begin{figure}[t]
\centering
% --- placeholder figure ---
\fbox{\parbox[c][6.0cm][c]{0.95\linewidth}{\centering
\textit{Figure placeholder.}\\
Pareto frontier in the (NDCG@10, weighted spread) plane on Ciao.\\
\texttt{[run\_outputs/ciao\_full/ciao\_pareto\_frontier.svg]}
}}
\caption{Pareto frontier on Ciao. Joint methods (JRSO,
JRSO\_second) dominate the line connecting the single-objective
baselines.}
\label{fig:pareto}
\end{figure}

\subsection{Training dynamics}
Figure~\ref{fig:curves} shows the per-epoch recommendation loss,
spread bonus and validation NDCG@10 for the four trained models. The
joint variants converge to a similar recommendation loss as
\textsc{RecSys-Only} but exhibit a clearly increasing spread-bonus
trajectory, confirming that the spread head learns a meaningful
gradient signal once the recommender warms up
($t\!\ge\!T_w$).

\begin{figure}[t]
\centering
% --- placeholder figure ---
\fbox{\parbox[c][6.0cm][c]{0.95\linewidth}{\centering
\textit{Figure placeholder.}\\
Training curves: recommendation loss (left), spread bonus (middle),
validation NDCG@10 (right) for all four trained models.\\
\texttt{[run\_outputs/ciao\_full/ciao\_training\_curves.svg]}
}}
\caption{Training curves on Ciao.}
\label{fig:curves}
\end{figure}

\subsection{Recommendation comparison bar chart}
Figure~\ref{fig:reccmp} reports NDCG@10 across all six methods. The
joint methods are within $1.5\!-\!8\%$ of the
\textsc{RecSys-Only} ceiling on NDCG@10 (a small recommendation cost)
while delivering substantial spread improvements.

\begin{figure}[t]
\centering
% --- placeholder figure ---
\fbox{\parbox[c][4.5cm][c]{0.95\linewidth}{\centering
\textit{Figure placeholder.}\\
NDCG@10 across methods.\\
\texttt{[run\_outputs/ciao\_full/ciao\_recommendation\_comparison.svg]}
}}
\caption{Recommendation comparison.}
\label{fig:reccmp}
\end{figure}

\subsection{Influence comparison bar chart}
Figure~\ref{fig:imcmp} reports weighted spread rate across all methods.
\textsc{JRSO} achieves the highest weighted spread among trainable
methods, exceeding even \textsc{IM-Only} on this metric while keeping
non-trivial recommendation quality.

\begin{figure}[t]
\centering
% --- placeholder figure ---
\fbox{\parbox[c][4.5cm][c]{0.95\linewidth}{\centering
\textit{Figure placeholder.}\\
Weighted spread rate across methods.\\
\texttt{[run\_outputs/ciao\_full/ciao\_influence\_comparison.svg]}
}}
\caption{Influence comparison.}
\label{fig:imcmp}
\end{figure}

\subsection{Statistical significance}
Table~\ref{tab:sig} reports paired permutation tests
($N\!=\!2{,}000$) on the per-item weighted-spread samples. All four
JRSO comparisons have $p\!<\!0.05$, confirming that the spread
improvements over the strongest single-objective baseline are not due
to chance.

\begin{table}[t]
\centering
\caption{Paired permutation tests on per-item weighted spread.
Significance is at $p<0.05$.}
\label{tab:sig}
\begin{tabular}{lcc}
\toprule
\textbf{Comparison} & $\Delta$ & $p$-value \\
\midrule
\textsc{JRSO\_second} vs \textsc{RecSys-Only} & $+30.51$ & $0.0075$ \\
\textsc{JRSO\_second} vs \textsc{IM-Only}     & $+9.30$  & $0.5052$ \\
\textsc{JRSO}         vs \textsc{RecSys-Only} & $+34.72$ & $0.0100$ \\
\textsc{JRSO\_third}  vs \textsc{Two-Stage}   & $+35.83$ & $0.0080$ \\
\bottomrule
\end{tabular}
\end{table}

\subsection{Ablation: edge subsampling}
Training-time edge subsampling reduces the per-step cost of the
differentiable IC diffusion from $\mathcal{O}(|E|)$ to
$\mathcal{O}(M_d)$. On \texttt{ciao\_full} ($|E|\!=\!146{,}800$,
$M_d\!=\!20{,}000$) this yields a $\sim\!7\times$ speed-up on a single
training step with no measurable degradation in NDCG@10 or weighted
spread (within run-to-run noise). At evaluation time we always use the
full graph.

\subsection{Discussion}
The experiments validate the central thesis of the paper: a model that
is jointly trained for recommendation and seeding is strictly more
useful than two separately optimised models. The observed price for
this joint training is a $1\!-\!8\%$ NDCG@10 cost relative to a pure
recommender, which is more than offset by a $25\!-\!30\%$ improvement
in weighted spread. Crucially, JRSO requires only one round of training
versus Two-Stage's two rounds and matches or beats Two-Stage on the
combined score.

The DiffNet++-style fusion in \textsc{JRSO\_third} is most useful when
the trust graph is dense relative to the rating matrix (a regime in
which the social signal carries non-trivial collaborative-filtering
information). On the Ciao smoke configuration the trust graph is
relatively sparse, which is why \textsc{JRSO\_second} dominates
\textsc{JRSO\_third} in this configuration. We expect the relative
ordering of \textsc{JRSO\_third} to improve on the larger
\texttt{ciao\_full} configuration, which retains roughly $2\times$ more
trust edges per user.

\subsection{Limitations and future work}
Three directions are particularly promising:
\begin{enumerate}
    \item Replacing the one-hop coverage in \eqref{eq:greedy} with a
    differentiable submodular surrogate so the discrete inference-time
    seed selection is itself end-to-end trained.
    \item Extending the IC relaxation to the Linear-Threshold
    family~\cite{shakarian2015independent}, which would make JRSO
    applicable to long-horizon non-progressive diffusion (epidemic-style
    spreading).
    \item Combining the joint loss with a multi-armed bandit
    explorer~\cite{IM_bandits,IM_bandits2,IM_bandits3} so that, in
    deployment, the model can refine its $\widetilde P$ estimates from
    observed downstream adoptions.
\end{enumerate}

% --------------------------------------------------------------------
\section{Conclusion}
% --------------------------------------------------------------------
We presented JRSO, an end-to-end framework that unifies socially-aware
recommendation and influence maximisation through a continuous seeding
relaxation, a differentiable Independent Cascade (Message-Passing
Diffusion) layer, and a LightGCN-style bipartite social encoder whose
weights also parameterise the IC edge probabilities. On the Ciao
trust/rating benchmark, JRSO and its DiffNet++-inspired variants
simultaneously match or improve recommendation quality and weighted
influence spread relative to single-objective baselines and a strong
two-stage pipeline, with statistically significant gains on weighted
spread. The framework is data-driven, treats the diffusion process as
endogenous to the recommender, and is the first system, to our
knowledge, to optimise both objectives within a single back-propagation.

\bibliographystyle{IEEEtran}
\bibliography{references}

\end{document}
